\documentclass[12pt,a4paper]{article}
\usepackage{fontspec}
\setmainfont{Times New Roman}
\usepackage[top=2.0cm,bottom=2.0cm,left=2.5cm,right=2.0cm]{geometry}
\usepackage{enumitem}
\usepackage{tabularx}
\usepackage{xltabular}
\usepackage{booktabs}
\usepackage{array}
\newcolumntype{Y}{>{\centering\arraybackslash}X}
\newcolumntype{C}[1]{>{\centering\arraybackslash}p{#1}}
\newcolumntype{L}[1]{>{\raggedright\arraybackslash}p{#1}}
\newcolumntype{R}[1]{>{\raggedleft\arraybackslash}p{#1}}
\pagestyle{plain}
\linespread{1.15}
\setlength{\parskip}{4pt plus 1pt minus 1pt}
\sloppy

\begin{document}
% --- HEADER 2 CỘT CHUẨN NGHỊ ĐỊNH 30/2020/NĐ-CP ---
\noindent
\begin{minipage}[t]{0.46\textwidth}
    \begin{center}
        \fontsize{11pt}{13pt}\selectfont
        CÔNG AN THÀNH PHỐ HÀ NỘI\\
        \textbf{PHÒNG CẢNH SÁT HÌNH SỰ (PC02)}\\[-0.2em]
        \rule{3.2cm}{0.6pt}\\[0.25cm]
        \fontsize{11pt}{13pt}\selectfont Số: \textbf{00/HS}
    \end{center}
\end{minipage}
\hfill
\begin{minipage}[t]{0.52\textwidth}
    \begin{center}
        \fontsize{11pt}{13pt}\selectfont
        \textbf{CỘNG HÒA XÃ HỘI CHỦ NGHĨA VIỆT NAM}\\
        \textbf{Độc lập - Tự do - Hạnh phúc}\\[-0.2em]
        \rule{3.6cm}{0.6pt}\\[0.25cm]
        \textit{Hà Nội, ngày 25 tháng 07 năm 2016}
    \end{center}
\end{minipage}

\vspace{0.5cm}

\begin{center}
    \textbf{\Large BIÊN BẢN LẤY LỜI KHAI}\\[0.15cm]
    \textit{(Lần thứ 02)}
\end{center}

\vspace{0.3cm}

Vào hồi 09 giờ 00 phút, ngày 26 tháng 07 năm 2016.

Tại: Phòng Cảnh sát Hình sự — Công an TP. Hà Nội.


\textbf{Chúng tôi gồm:}

\begin{itemize}[leftmargin=*, itemsep=2pt, topsep=2pt, parsep=0pt]
  \item Điều tra viên: \textbf{Lê Minh} — Chức vụ: Đại úy, Điều tra viên Đội 5 (Đội Điều tra Trọng án).
  \item Cán bộ ghi biên bản: \textbf{Nguyễn Văn Hoàng} — Chức vụ: Trung úy, Cán bộ Đội 5.

\end{itemize}
\textbf{Tiến hành lấy lời khai của người có quyền lợi, nghĩa vụ liên quan:}

\begin{itemize}[leftmargin=*, itemsep=2pt, topsep=2pt, parsep=0pt]
  \item Họ và tên: \textbf{TRẦN THỊ HÀ} | Giới tính: Nữ
  \item Sinh ngày: 22/09/1990 tại Hà Nội.
  \item Quốc tịch: Việt Nam | Dân tộc: Kinh | Tôn giáo: Không.
  \item CMND số: \textbf{0906438125} do CA TP. Hà Nội cấp ngày 10/05/2006.
  \item Nơi ĐKHKTT: Số 8, Ngõ 12, Đường Bờ Kè, Phường Phân khu Cảng, Quận Sông Hồng, TP. Hà Nội.
  \item Chỗ ở hiện tại: Số 8, Ngõ 12, Đường Bờ Kè, Phường Phân khu Cảng, Quận Sông Hồng, TP. Hà Nội.
  \item Nghề nghiệp: Kế toán tổng hợp (Công ty TNHH Vận tải Sông Hồng).
  \item Mối quan hệ với bị hại: Bạn gái của bị hại (Nguyễn Văn Khang).


\end{itemize}

\vspace{0.3cm}\noindent\textbf{\large NỘI DUNG LỜI KHAI (HỎI VÀ ĐÁP)}\\[0.15cm]


\noindent\textbf{Hỏi (ĐTV Lê Minh):} \textit{Chị Hà, trong lần khai báo đầu tiên, chị khẳng định tối 24/07 chị ở trong phòng trọ xem phim bộ trên VTV3 rồi đi ngủ. Chị có muốn bổ sung hay thay đổi nội dung lời khai nào không?}\\[0.15cm]

\noindent\textbf{Đáp (Trần Thị Hà):} \textit{(Bình tĩnh, giọng nhẹ nhàng)} Thưa cán bộ, em đã khai đúng sự thật rồi ạ. Tối đó em ở phòng, xem phim rồi đi ngủ.\\[0.35cm]


\noindent\textbf{Hỏi (ĐTV Lê Minh):} \textit{Lúc 20 giờ 32 phút tối hôm đó, chị có gọi điện thoại hoặc gửi tin nhắn cho nạn nhân Khang không?}\\[0.15cm]

\noindent\textbf{Đáp (Trần Thị Hà):} \textit{(Hơi ngập ngừng)} ...Có ạ. Em nhớ ra rồi, tối đó em có gọi cho anh Khang một cuộc nhưng anh ấy không nghe máy. Máy tự động chuyển sang hộp thư thoại... Em nhớ là em có để lại một đoạn tin nhắn thoại ngắn hỏi anh ấy đang làm gì mà không nghe máy của em.\\[0.35cm]


\noindent\textbf{Hỏi (ĐTV Lê Minh):} \textit{Cơ quan điều tra đã trưng cầu Phòng Kỹ thuật Hình sự giám định đoạn tin nhắn thoại chị để lại lúc 20 giờ 32 phút. Kết quả phân tích âm thanh cho thấy trong đoạn ghi âm có lẫn tiếng còi tàu hỏa và tiếng chuông cảnh báo rào chắn đường sắt. Đối chiếu với lịch trình đoàn tàu chạy qua giao cắt Đường Bờ Sông lúc 20 giờ 30 phút và kết quả đo cường độ âm thanh thực địa, tiếng còi tàu trong ghi âm (68 dB) chỉ có thể ghi nhận tại vị trí trước cổng nhà ông Khang (cách đường ray 30m), chứ không thể thu từ phòng trọ của chị cách đường ray 1,2 km. Chị giải thích thế nào về việc này?}\\[0.15cm]

\noindent\textbf{Đáp (Trần Thị Hà):} \textit{(Mặt tái đi, tay bắt đầu run, mắt nhìn xuống sàn)} ...Cán bộ... cán bộ nói gì em không hiểu... Phòng trọ em ở gần đường ray mà... lâu lâu em vẫn nghe thấy tiếng tàu...\\[0.35cm]


\noindent\textbf{Hỏi (ĐTV Lê Minh):} \textit{Theo hồ sơ xác minh, phòng trọ của chị tại số 8 Ngõ 12 Đường Bờ Kè cách đường ray tàu hỏa 1,2 kilômét. Tổ công tác đã đo đạc thực tế: tiếng còi tàu tại vị trí phòng trọ chỉ còn dưới 25 dB là không đủ lọt vào micro điện thoại. Trong khi cường độ thu được trong ghi âm là 68 dB, tương đương vị trí đứng cách đường ray khoảng 30 mét, chính xác là trước cổng nhà nạn nhân. Chị có thể giải thích sự chênh lệch này không?}\\[0.15cm]

\noindent\textbf{Đáp (Trần Thị Hà):} \textit{(Hai tay vặn chặt vạt áo, giọng bắt đầu lắp bắp)} Em... em... không biết! Có khi máy ghi sai... Hoặc gió mang tiếng tàu đi xa hơn bình thường... Em khẳng định tối đó em chỉ ở trong phòng chứ không đi đâu hết!\\[0.35cm]


\noindent\textbf{Hỏi (ĐTV Lê Minh):} \textit{Thêm một vấn đề nữa. Chị khai rằng tối 24/07, chị xem phim bộ trên kênh VTV3. Tuy nhiên, lịch phát sóng chính thức của VTV3 tối Thứ Sáu ngày 24/07/2016 là chương trình giải trí, không có phim bộ nào. Chị giải thích thế nào?}\\[0.15cm]

\noindent\textbf{Đáp (Trần Thị Hà):} \textit{(Giật mình, mắt mở to, mặt đỏ bừng)} Em... em nhớ nhầm kênh... Có thể em xem kênh khác... Em không nhớ rõ, tối đó em buồn lắm, em nhớ lung tung hết...\\[0.35cm]


\noindent\textbf{Hỏi (ĐTV Lê Minh):} \textit{Tức là chị thừa nhận lời khai lần 1 về việc xem phim bộ VTV3 là không chính xác?}\\[0.15cm]

\noindent\textbf{Đáp (Trần Thị Hà):} \textit{(Im lặng kéo dài, hai vai run rẩy. Bất ngờ ngẩng lên, giọng gắt gỏng)} Em nhầm kênh thôi chứ em có nói dối đâu! Cán bộ hỏi gì mà cứ dồn ép em! Em đã bảo tối đó em ở phòng trọ mà! \textit{(Gạt nước mắt)}\\[0.35cm]


\noindent\textbf{Hỏi (ĐTV Lê Minh):} \textit{Không ai dồn ép chị cả. Chúng tôi chỉ muốn làm rõ sự thật. Tôi hỏi lại lần cuối: Tối 24/07/2016, trong khoảng thời gian từ 19 giờ đến 21 giờ, chị có đến khu vực nhà nạn nhân Nguyễn Văn Khang tại số 14 Đường Bờ Sông không?}\\[0.15cm]

\noindent\textbf{Đáp (Trần Thị Hà):} \textit{(Đập tay xuống bàn, mắt đỏ ngầu, giọng nghẹn ngào)} EM KHÔNG ĐẾN ĐÓ! Em không biết gì hết! Em không muốn trả lời nữa! Cán bộ muốn bắt em thì cứ bắt, em không khai gì thêm!\\[0.35cm]

\textit{(Đối tượng Trần Thị Hà từ chối trả lời mọi câu hỏi tiếp theo. Tay vặn vẹo vạt áo liên tục, mắt nhìn chằm chằm xuống mặt bàn, thi thoảng rung vai khóc nấc.)}



Buổi lấy lời khai kết thúc vào hồi 09 giờ 45 phút cùng ngày. Biên bản này đã được đọc lại cho chị Trần Thị Hà và những người tham dự cùng nghe, công nhận ghi đúng lời khai.


Chị Trần Thị Hà từ chối ký tên vào biên bản với lý do: \textit{"Tôi không muốn ký tên vào biên bản này"}.


\textbf{NGƯỜI CHỨNG KIẾN XÁC NHẬN VIỆC TỪ CHỐI KÝ}

\textit{(Ký, ghi rõ họ tên)}

\textit{(Xác nhận chị Trần Thị Hà đã nghe đọc biên bản nhưng từ chối ký tên)}



\textit{(Đã ký)}

\textbf{Phạm Văn Đức}

\textit{(Cán bộ trực ban Trụ sở)}

\vspace{0.8cm}
\noindent
\begin{minipage}[t]{0.48\textwidth}
    \begin{center}
        \textbf{CÁN BỘ GHI BIÊN BẢN}\\[0.1cm]
        \textit{(Ký, ghi rõ họ tên)}\\[1.6cm]
        \textbf{Nguyễn Văn Hoàng}
    \end{center}
\end{minipage}
\hfill
\begin{minipage}[t]{0.48\textwidth}
    \begin{center}
        \textbf{ĐIỀU TRA VIÊN}\\[0.1cm]
        \textit{(Ký, ghi rõ họ tên)}\\[1.6cm]
        \textbf{Lê Minh}
    \end{center}
\end{minipage}
\end{document}
